\section{Annexes}

\subsection{Annexe A : Constante de raideur}

\begin{table}[H]
    \centering
    \includegraphics[scale=0.6]{tab1.png}
	\captionof{table}{Mesures pour determiner constante de raideur k}
\end{table}

\subsection{Annexe B : Libres}

\begin{minipage}[t]{0.46\textwidth}
    \centering
    \includegraphics[width=\linewidth]{libre_frein_0_1A.png}
        \captionof{figure}{Mesure de l'angle theta a 0.1A}
	\label{fig:oscillations}
\end{minipage}
\hfill
\begin{minipage}[t]{0.46\textwidth}
    \centering
    \includegraphics[width=\linewidth]{libre_frein_0_2A.png}
        \captionof{figure}{Mesure de l'angle theta a 0.2A}
\end{minipage}

\begin{table}[!ht]
\centering
\footnotesize
\setlength{\tabcolsep}{4pt}
\renewcommand{\arraystretch}{0.9}
\begin{tabular}{|c|c|c|c|c|}
\hline
\textbf{Fichier} 
& \textbf{$\omega_0$ [rad/s]} 
& \textbf{$\sigma_{\omega_0}$ [rad/s]} 
& \textbf{$f_0$ [Hz]} 
& \textbf{$\sigma_{f_0}$ [Hz]} \\
\hline
0.1A & 3.25153 & 0.00004 & 0.51745 & 0.000006 \\
0.2A & 3.25250 & 0.00006 & 0.51760 & 0.000010 \\
0.3A & 3.25089 & 0.00009 & 0.51734 & 0.000014 \\
0.4A & 3.25255 & 0.00019 & 0.51760 & 0.000030 \\
0.5A & 3.25308 & 0.00028 & 0.51768 & 0.000045 \\
0.9A & 3.12291 & 0.01293 & 0.49734 & 0.00206 \\
\hline
\textbf{Moyenne} 
& \textbf{3.25} 
& \textbf{0.02} 
& \textbf{0.517} 
& \textbf{0.003} \\
\hline
\end{tabular}
\caption{Valeurs calculées de $\omega_0$ et $f_0$ à partir de $\omega$ et $\lambda$.}
\label{val:f0}
\end{table}


\begin{thebibliography}{99}

\bibitem{schemapohl}
Gilbert Gastebois, \textit{Le pendule de Pohl},
\\
                \texttt{https://ggastebois.fr/java/pendpohl/theorie\_pend\_pohl.pdf}
\\
        Consulte le : 17 Novembre 2025

\end{thebibliography}

